\documentclass[10pt,a4paper]{amsart}
\usepackage[margin=19mm]{geometry}
\usepackage{amsmath,amssymb,mathtools}
\usepackage[scr=rsfs,cal=euler]{mathalfa}
\usepackage{xfrac}
\usepackage[unicode,hidelinks,bookmarks=false]{hyperref}
\newcommand{\ie}{i.e.\ }
\newcommand{\cf}{cf.\ }
\newcommand{\resp}{respectively\ }
\newcommand{\Subsection}{\subsection}
\providecommand{\nobreakdash}{\nobreak\hbox{-}}
\setlength{\emergencystretch}{2em}
\multlinegap0pt
\usepackage{bm}
\usepackage{bbm}
\usepackage{dsfont}
\usepackage{xspace}
\usepackage{cancel}
\usepackage{mathtools}
\usepackage[shortlabels]{enumitem}
\usepackage{amsthm}
\makeatletter
\@ifundefined{lemma}{\newtheorem{lemma}{Lemma}}{}
\@ifundefined{theorem}{\newtheorem{theorem}[lemma]{Theorem}}{}
\makeatother
\newcommand{\B}{{\mathcal{B}}}
\newcommand{\bC}{{\mathbb{C}}}
\newcommand{\bM}{{\mathbb{M}}}
\newcommand{\rC}{\mathrm{C}}
\renewcommand{\S}{{\mathcal{S}}}
\renewcommand{\phi}{\varphi}
\title[A new obstruction to Arveson's hyperrigidity conjecture]{A new obstruction to Arveson's hyperrigidity conjecture}
\author{Rapha\"el Clou\^atre}
\date{Source: arXiv:2509.19238v1 (CC BY 4.0)}
\begin{document}
\maketitle

\noindent\emph{Source and licence.} A new obstruction to Arveson's hyperrigidity conjecture. Version arXiv:2509.19238v1 (CC BY 4.0), distributed under the Creative Commons Attribution 4.0 International licence, as recorded in the arXiv licence metadata for this exact version and shown on the abstract page https://arxiv.org/abs/2509.19238v1. Full rights notice: licenses/arxiv-2509.19238-CC-BY-4.0.txt.\par\smallskip

\noindent\textit{Editorial scope.} Counted unit: the Theorem whose source label is \texttt{T:singrepuep} in the article (source0.tex lines 620--626, source label \texttt{T:singrepuep}), delivered together with the article's own complete original proof of it (source0.tex lines 627--733, one \texttt{proof} environment of 107 lines). No mathematics is written, repaired or re-derived: statement and proof are the article's own text line for line. Required context retained uncounted: L:JinB (lines 359--361); Eq:adapted (lines 527--531); L:nonunital (lines 596--601). Every definition, display and label of those lines is retained unchanged. Omitted from this package and not redistributed: 1--358, 362--526, 532--595, 602--619, 734--1023. Nothing inside a retained excerpt is omitted or altered: every retained body is delivered exactly as the article's lines are written, so no omission receipt of any kind accompanies this package. Citations: the retained excerpts carry no citation command at all, so no bibliography entry of the article is transcribed and no reference is redistributed. Reference adaptation: none was needed and none was made; every reference inside the delivered unit resolves inside it, so no unresolved reference or citation is produced. Printed slips kept literal: no printed word, symbol, hypothesis or number of the counted statement or of its proof was altered. Layout and class: the article's own class is \texttt{amsart} with packages including graphicx, amscd, amsmath, amsthm, amsfonts, amssymb, mathrsfs, bm, enumerate, amsrefs, xcolor, hyperref, inputenc, tikz; the wrapper is the lane's pinned \texttt{amsart} editorial wrapper at 10pt on A4, which restates the theorem-like environments the retained text needs and adds the article's own macro definitions that the retained text needs (\texttt{\textbackslash B}, \texttt{\textbackslash S}, \texttt{\textbackslash bC}, \texttt{\textbackslash bM}, \texttt{\textbackslash phi}, \texttt{\textbackslash rC}), with extra wrapper packages (bm, bbm, dsfont, xspace, cancel, mathtools, enumitem). The delivered numbering is the unit's own. Assets: no retained span contains a figure environment, a graphics-inclusion command, a PDF-page inclusion, a table or a TikZ picture; no image file is copied into this package, the package declares no asset dependency and no image file is redistributed with it. Licence: version arXiv:2509.19238v1 of this article is distributed under the Creative Commons Attribution 4.0 International licence, as recorded in the arXiv licence metadata for this exact version and shown on the abstract page https://arxiv.org/abs/2509.19238v1. A full rights notice is delivered at licenses/arxiv-2509.19238-CC-BY-4.0.txt. Characters: every retained excerpt is pure ASCII, so the delivered engine, XeLaTeX with its default Latin Modern text font, reports no missing character.
\begin{lemma}\label{L:JinB}
The $\rC^*$-algebra $\B(X,J)$ consists of elements of the form $\begin{bmatrix} a & b \\ c & d\end{bmatrix}$ with $a\in \rC^*(X)+J$, $b\in J, c\in J$ and $d\in J+\bC 1$.
\end{lemma}

\begin{equation}\label{Eq:adapted}
\Pi\left(\begin{bmatrix}
a& b \\ c & d
\end{bmatrix}\right)-\chi(d)1=\pi(a-\chi(d)1)
\end{equation}

\begin{lemma}\label{L:nonunital}
Let $A$ be a unital $\rC^*$-algebra and let $\phi:A\to B(H)$ be a contractive completely positive map. Let $e=\phi(1)$. Then,
\[
\phi(a)=\phi(a)e=e\phi(a), \quad a\in A.
\]
\end{lemma}

\begin{theorem}\label{T:singrepuep}
Let $A$ be a unital $\rC^*$-algebra containing a closed two-sided ideal $J$ and an operator system $X$. Let $\Pi:\B(X,H)\to B(H)$ be a unital $*$-representation annihilating $\bM_2(J)$, and let $\pi:\rC^*(X)+J\to B(H)$ be the unique $*$-representation such that the pair $(\Pi,\pi)$ is $(X,J)$-adapted. Then, the following statements are equivalent.
\begin{enumerate}[{\rm (i)}]
\item $\Pi$ has the unique extension property with respect to $\S(X,J)$.
\item $\pi$ is the  unique contractive completely positive extension of $\pi|_X$ to $\rC^*(X)+J$.
\end{enumerate}
\end{theorem}

\begin{proof}
(ii)$\Rightarrow$(i): Let $\Phi:\B(X,J)\to B(H)$ be a unital completely positive map agreeing with $\Pi$ on $\S(X,J)$.
By Lemma \ref{L:JinB}, we may define a contractive completely positive map $\phi:\rC^*(X)+J\to B(H)$ as
\[
\phi(a)=\Phi\left(\begin{bmatrix}
a& 0 \\ 0 & 0
\end{bmatrix}\right), \quad a\in \rC^*(X)+J.
\]
Since the pair $(\Pi,\pi)$ is $(X,J)$-adapted,  \eqref{Eq:adapted} implies that $\phi$ agrees with $\pi$ on $X$, so our assumption forces $\phi$ and $\pi$ to agree on $\rC^*(X)+J$.



 Fix $\begin{bmatrix}
a & b \\ c& d
\end{bmatrix}\in \B(X,J)$. By Lemma \ref{L:JinB}, we infer that $a\in \rC^*(X)+J$, $b\in J,c\in J$ and $d\in J+\bC 1$. Thus, the previous paragraph implies that
\begin{equation}\label{Eq:piphia}
\pi(a-\chi(d)1)=\phi(a-\chi(d)1).
\end{equation}
Moreover the elements
$\begin{bmatrix}
0& 0 \\ 0 & d-\chi(d)1
\end{bmatrix},\begin{bmatrix}
0 & 0 \\ 0 & b^*b
\end{bmatrix},\begin{bmatrix}
0 & 0 \\ 0 & cc^*
\end{bmatrix}$ all lie in $ \bM_2(J)\cap \S(X,J)$
so
\[
\Phi\left(\begin{bmatrix}
0 & 0 \\  0& d-\chi(d)1
\end{bmatrix}\right)=\Phi\left(\begin{bmatrix}
0 & 0 \\ 0 & b^*b
\end{bmatrix} \right)=\Phi\left(\begin{bmatrix}
0 & 0 \\ 0 & cc^*
\end{bmatrix} \right)=0
\]
since  $\Phi$ and $\Pi$ agree on $\S(X,J)$ and $\Pi$  annihilates $\bM_2(J)$.
By the Schwarz inequality, this implies that
\[
\Phi\left(\begin{bmatrix}
0& b \\ 0& 0
\end{bmatrix}\right)=\Phi\left(\begin{bmatrix}
0& 0 \\ c& 0
\end{bmatrix}\right)=0
\]
whence
\begin{equation}\label{Eq:phi}
\Phi\left(\begin{bmatrix}
a& b \\ c & d
\end{bmatrix}\right)=\Phi\left(\begin{bmatrix}
a& 0 \\ 0 & \chi(d)1
\end{bmatrix}\right)=\Phi\left(\begin{bmatrix}
a-\chi(d)1& 0 \\ 0 & 0
\end{bmatrix}\right)+\chi(d)1.
\end{equation}
Employing Equations \eqref{Eq:adapted},\eqref{Eq:piphia} and \eqref{Eq:phi} we find
\begin{align*}
\Pi\left(\begin{bmatrix}
a& b \\ c & d
\end{bmatrix}\right)-\chi(d) 1&=\pi(a-\chi(d)1)=\phi(a-\chi(d)1)\\
&=\Phi\left(\begin{bmatrix}
a-\chi(d)1& 0 \\ 0 & 0
\end{bmatrix}\right)\\
&=\Phi\left(\begin{bmatrix}
a& b \\ c & d
\end{bmatrix}\right)-\chi(d)1.
\end{align*}
We conclude that $\Pi$ and $\Phi$ agree on $\B(X,J)$, so indeed $\Pi$ has the unique extension property with respect to $\S(X,J)$.


(i)$\Rightarrow$(ii):   Let $\phi:\rC^*(X)+J \to B(H)$ be a contractive completely positive map agreeing with $\pi$ on $X$. By Lemma \ref{L:nonunital}, the projection $p=\pi(1)$ commutes with the image of $\phi$.
Invoking Lemma \ref{L:JinB}, we can then define a \emph{unital} completely positive map $\Phi:\B(X,J)\to B(H)$ as
\begin{equation}\label{Eq:phi0}
\Phi\left(\begin{bmatrix}
a& b \\ c & d
\end{bmatrix}\right)=\phi(a)p+\chi(d)(1-p)
\end{equation}
 for $\begin{bmatrix}
a& b \\ c & d
\end{bmatrix}\in \B(X,J)$.
Using \eqref{Eq:adapted} we may also write
\begin{align*}
\Pi\left( \begin{bmatrix}
a& b \\ c & d
\end{bmatrix}\right)
%&=\Pi\left(\begin{bmatrix}
%a& 0 \\ 0 & \chi(d)1
%\end{bmatrix}\right)\\
%&=\pi(a)+\chi(d)\left(1-\Pi\left(\begin{bmatrix}
%1& 0 \\ 0 &0
%\end{bmatrix}\right)\right)\\
&=\pi(a)p+\chi(d)(1-p).
\end{align*}
It then follows from \eqref{Eq:phi0} that $\Phi$ and $\Pi$ agree on $\S(X,J)$, so by assumption $\Phi$ and $\Pi$ agree on $\B(X,J)$. 
Finally, for $a\in \rC^*(X)+J$ we have $\begin{bmatrix}
a& 0 \\ 0 & 0
\end{bmatrix}\in  \B(X,J)$ by  Lemma \ref{L:JinB}, so an application of Lemma \ref{L:nonunital} and Equation \eqref{Eq:phi0} implies
\begin{align*}
\phi(a)&=\phi(a)p=\Phi\left(\begin{bmatrix}
a& 0 \\ 0 & 0
\end{bmatrix}\right)
=\Pi\left(\begin{bmatrix}
a& 0 \\ 0 & 0
\end{bmatrix}\right)=\pi(a).
\end{align*}
We conclude that indeed $\pi$ is the unique contractive completely positive extension of $\pi|_X$ to $\rC^*(X)+J$.
\end{proof}

\end{document}
